\begin{longtable}{@{}l p{12.5cm}@{}}
  \caption{Skills per-instance answer comparison ($\Delta \geq 75\%$) for the 1 most divergent items: same question, each model's answer and reasoning.}\label{tab:cmp_skills_pos75} \\
  \toprule
  \endfirsthead
  \toprule
  \endhead
  \multicolumn{2}{@{}l}{\textbf{94826} \quad $\Delta = +80\%$} \\
  \textbf{Question} & You are an evaluator using an essay rubric to assess a kindergarten student's short narrative (about 4--6 sentences) written in simple English. Analyze the student's essay by: 1) identifying the specific rubric criteria (e.g., organization, vocabulary, sentence structure, conventions, creativity) that are met and those that are not met; 2) citing at least three concrete pieces of textual evidence (short quoted phrases or sentences from the essay) to support each claim; and 3) explaining how each piece of evidence maps to the rubric descriptors (showing the reasoning that connects evidence to criterion). Provide a clear, organized analysis that could be shared with the classroom teacher to justify your scoring decisions. \\
  \textbf{Reference} & Student essay (example): "I went to the park. I saw a big dog. The dog ran fast. I felt happy and I played with the dog." \newline Analysis: \newline 1) Organization: Partially met. Evidence: "I went to the park." (opening setting), "I saw a big dog." (event follows setting), and "I felt happy and I played with the dog." (concluding feeling/action). These three phrases show a basic sequence: setting event reaction, which aligns with the rubric descriptor for simple narrative organization (clear beginning, middle, end). However, the transitions between ideas are minimal and there is no explicit ending sentence that wraps up the story, so the essay meets the organization criterion at a basic level but not at a more developed level. \newline 2) Vocabulary and Word Choice: Mostly met for age level. Evidence: "big dog,ran fast," and "happy." These concrete, age-appropriate words convey the main images and emotions clearly. The use of descriptive words like "big" and "fast" demonstrates an attempt to add detail. The vocabulary is simple but effective for kindergarten expectations, mapping to a rubric descriptor that values appropriate and concrete word choice for the grade. \newline 3) Sentence Structure and Fluency: Partially met. Evidence: "I went to the park." (complete simple sentence), "I saw a big dog." (complete simple sentence), and "I felt happy and I played with the dog." (compound sentence connected with "and"). The presence of complete sentences shows control of basic sentence formation; the compound sentence indicates emerging ability to combine ideas. However, all sentences are short and formulaic, so while the student demonstrates foundational sentence structure, they do not yet show varied or complex sentences required for higher rubric levels. \newline 4) Conventions (spelling, capitalization, punctuation): Mostly met. Evidence: Each sentence begins with a capital "I" and ends with a period in the sample text provided. Words like "park,dog," and "happy" are spelled correctly. This matches rubric descriptors for correct use of basic conventions at the kindergarten level. One possible area for attention is spacing and any capitalization of proper nouns if present in other drafts, but in this sample conventions are largely correct. \newline \textit{[\ldots\ truncated \ldots]} \newline - The absence of additional sensory details, characterization, or varied sentence types maps to rubric levels indicating minimal elaboration and limited complexity. \newline Overall evaluative synthesis: The essay demonstrates solid foundational skills expected of a kindergarten writer: clear simple sentences, appropriate basic vocabulary, correct basic conventions, and a recognizable narrative sequence (beginning event reaction). Areas for growth include adding more descriptive detail (who the dog was, what it looked like, how it acted) and varying sentence length/structure to increase elaboration and narrative interest. The textual evidence cited above supports awarding middle-level ratings for organization, vocabulary, and conventions, lower-middle ratings for sentence structure, and a low rating for creativity/elaboration. \\
  \midrule
  \textbf{Kimi-K2.6} & score 1/10 \\
  & [1-2] The student's answer completely omits the required analysis, evidence, and mapping, making the result meaningless for the evaluation task due to the missing input. \\
  \textbf{Innospark-1T} & score 9/10 \\
  & Let me think through this carefully. \newline The question/task: You are an evaluator using an essay rubric to assess a kindergarten student's short narrative (about 4--6 sentences) written in simple English. Analyze the student's essay by: 1) identifying the specific rubric criteria (e.g., organization, vocabulary, sentence structure, conventions, creativity) that are met and those that are not met; 2) citing at least three \newline My reasoning: Kindergarten Narrative Evaluation: Sample Analysis \newline Important Note \newline To provide a thorough analysis, I need a specific student essay to evaluate. Below is a sample analysis using a hypothetical kindergarten narrative. If you have a specific student essay you'd like me to evaluate, please provide it, and I'll analyze that work instead. Sample Student Essay (Hypothetical) \newline Title: "My Trip to the ZooI went to the zoo. I saw a big elephant. The elephant was gray. It had a long nose. I liked the elephant. We went home." \newline \textit{[\ldots\ truncated \ldots]} \newline --- \newline If you have a specific student essay you'd like evaluated using this same rigorous rubric-based approach, please share it, and I'll provide a detailed, evidence-based analysis tailored to that work. \\
  \specialrule{1.2pt}{8pt}{8pt}
  \bottomrule
\end{longtable}